% EDITA AQUÍ LA TIPOGRAFÍA Y LOS TAMAÑOS DE PORTADA.
% La misma familia se usa en texto, títulos y citas.
\newcommand{\FuenteTexto}{Latin Modern Sans}
\newcommand{\TamTituloPortada}{22}
\newcommand{\InterlineadoTituloPortada}{25}
\newcommand{\TamNombrePortada}{14}
\newcommand{\TamSubtituloPortada}{12}
\newcommand{\TamCursoPortada}{11}
\newcommand{\TamSemestrePortada}{10}
\newcommand{\NotaExpositor}{Todo lo dicho en esta sesión es responsabilidad del expositor.}

\setmainfont{\FuenteTexto}
\setsansfont{\FuenteTexto}
\renewcommand{\familydefault}{\sfdefault}
\usefonttheme{professionalfonts}
\setbeamerfont{normal text}{family=\sffamily,size=\normalsize}
\setbeamerfont{frametitle}{family=\sffamily,series=\bfseries}
\setbeamerfont{quote}{family=\sffamily,shape=\upshape,size=\normalsize}
\setbeamerfont{quotation}{family=\sffamily,shape=\upshape,size=\normalsize}

% Iconos de una biblioteca existente; no representan hechos históricos.
\usepackage{fontawesome5}

% Todos los rótulos son visibles: no heredan el color negro de los enlaces.
\setbeamertemplate{headline}{%
  \begin{beamercolorbox}[wd=\paperwidth,ht=3ex,dp=1ex]{section in head/foot}
    \sffamily\bfseries\hbox to\paperwidth{\hfil Motivación\hfil Práctica\hfil Interpretación\hfil Cierre\hfil}
  \end{beamercolorbox}
}
